%% kapa: exam=1 round=27 year=2016 period=2 form=A subject=accounting subject_ko=회계학 seq=29 number=69
%% source: Q-Net 감정평가사 기출문제 (cst00309) artlSeq=5169402 file=extracted/1008/5169402/제27회 감정평가사 1차 문제/제27회 감정평가사 2교시문제(A형).pdf
%% conversion: kordoc 4.24.51 -> tools/gen1.py (auto); review=proofread-sonnet; flags=none
\begin{question}{69}
충당부채와 우발부채에 관한 설명으로 옳지 않은 것은?
\begin{choices}
  \item 충당부채를 인식하기 위해서는 당해 의무를 이행하기 위하여 경제적효익을 갖는 자원이 유출될 가능성이 매우 높아야 한다.
  \item 우발부채는 경제적효익을 갖는 자원의 유출을 초래할 현재의무가 있는지의 여부가 아직 확인되지 아니한 잠재적 의무이므로 부채로 인식하지 않는다.
  \item 재무제표는 미래 시점의 예상 재무상태가 아니라 보고기간말의 재무상태를 표시하는 것이므로, 미래영업을 위하여 발생하게 될 원가에 대하여는 충당부채를 인식하지 않는다.
  \item 충당부채로 인식되기 위해서는 과거사건으로 인한 의무가 기업의 미래행위(즉, 미래 사업행위)와 독립적이어야 한다.
  \item 상업적 압력 때문에 공장에 특정 정화장치를 설치하기 위한 비용지출을 계획하고 있는 경우 공장운영방식을 바꾸는 등의 미래행위를 통하여 미래의 지출을 회피할 수 있으므로 당해 지출은 현재의무가 아니며 충당부채도 인식하지 아니한다.
\end{choices}
\end{question}
\answer{1}
